Bezüglich der Meereshöhe liegt der Gipfel um $h = \hevO$ höher:
\begin{align}
 E_\mathrm{pot} &= \EpF \\
   &= \mbs\cdot\gE\cdot\hev \\
   &= \Ep \approx \result{\EpP{6}{2}}
\end{align}
Bezüglich des Gipfels liegt die Meereshöhe um $h$ tiefer, die Höhe ist also $-h$ und die potentielle Energie
\begin{align}
 E_\mathrm{pot} &= -m\,g\,h = \Epn \approx \result{\EpnP{6}{2}}
\end{align}
Die potentielle Energie hängt vom gewählten Nullniveau ab; nur Differenzen der potentiellen Energie haben eine physikalische Bedeutung.
